%% chapters/05-hasil-pembahasan.tex
%% Bab V — Hasil dan Pembahasan
%% Aturan: lihat .cursor/rules/10-bab5-hasil-pembahasan.mdc

\chapter{Hasil dan Pembahasan}
\label{bab:hasil}

Bab ini menyajikan hasil eksperimen yang dirancang pada
Subbab~\ref{subsec:bab4_eksperimen}. Setiap eksperimen dilaporkan
dalam sub-subbab tersendiri yang dilengkapi dengan tabel atau gambar
beserta diskusi naratif yang mengaitkan temuan dengan pertanyaan
penelitian dan hipotesis.

\section{Hasil Eksperimen Utama}
\label{subsec:bab5_main}

\subsection{Perbandingan dengan SOTA — XJTU-SY}
\label{subsec:bab5_main_xjtusy}

% Tabel V.1: Perbandingan kinerja Mamba-xLSTM-Net vs baseline pada XJTU-SY.
% Diskusi: pembacaan tabel, pembanding terkuat, hubungan ke H1.
\dots

\subsection{Perbandingan dengan SOTA — PHM2012}
\label{subsec:bab5_main_phm2012}
\dots

\subsection{Visualisasi Kurva Prediksi RUL}
\label{subsec:bab5_kurva}
\dots

\section{Studi Ablasi}
\label{subsec:bab5_ablasi}

\subsection{Ablasi Komponen Arsitektur}
\label{subsec:bab5_ablasi_arch}
% Tabel ablasi: -Mamba, -mLSTM, -sLSTM, -fusi, dst.
\dots

\subsection{Ablasi Hyperparameter}
\label{subsec:bab5_ablasi_hp}
% Plot tren metrik vs hyperparameter (dimensi state, jumlah head, k SAE)
\dots

\section{Analisis Interpretabilitas}
\label{subsec:bab5_interp}

\subsection{Atribusi Fitur SHAP dan Integrated Gradients}
\label{subsec:bab5_shap}
% Bar plot atribusi global + plot temporal early/mid/late life
\dots

\subsection{Sparse Autoencoder: Konsep yang Ditemukan}
\label{subsec:bab5_sae_concepts}
% Distribusi sparsity, maximum-activating examples, interpretasi 3-5 fitur SAE
\dots

\subsection{Korespondensi dengan Frekuensi Karakteristik Bantalan}
\label{subsec:bab5_korespondensi}
% Spektrum + garis vertikal BPFO/BPFI/BSF teoretis + heatmap aktivasi SAE
\dots

\section{Generalisasi Lintas Dataset}
\label{subsec:bab5_generalisasi}
% Train PHM2012 -> test XJTU-SY, dan sebaliknya
\dots

\section{Pembahasan Lebih Lanjut}
\label{subsec:bab5_diskusi}

\subsection{Kekuatan dan Kelemahan Metode Usulan}
\label{subsec:bab5_pro_kontra}
\dots

\subsection{Posisi terhadap \emph{State-of-the-Art}}
\label{subsec:bab5_posisi}
\dots

\subsection{Implikasi Praktis}
\label{subsec:bab5_implikasi}
\dots

\section{Keterbatasan Penelitian}
\label{subsec:bab5_keterbatasan}

% Daftar jujur keterbatasan yang akan menjadi seed untuk Saran di Bab VI:
% - Validasi terbatas pada bantalan gelinding deep groove ball;
% - Asumsi degradasi monoton;
% - Kondisi operasi laboratorium;
% - Klaim interpretabilitas berbasis korelasi statistik (bukan kausal);
% - Komputasi pelatihan masih relatif mahal.
\dots
